\section*{Chapter \ref{ch:rc:build-a-risk-engine} --- Build: A Risk Engine}

\begin{solution}{exo:rc:build-a-risk-engine:1}
Full: 251 calls for either (the base and 250 scenarios). Grid: $1+8\times6=49$ for the swap (eight \omterm{def:rc:curve-construction:pillar}{pillars},
six non-zero points each) and $1+9\times6=55$ for the option (the spot too). Delta--gamma: $1+2\times8=17$ and
$1+2\times9=19$. Book: $1\,000\times251=251\,000$; $800\times49+200\times55=50\,200$; $800\times17+200\times19=17\,400$.
\end{solution}

\begin{solution}{exo:rc:build-a-risk-engine:2}
The grid prices each trade at a fixed set of shifts; each scenario then costs an interpolation and a sum,
not a pricing call. Adding scenarios adds arithmetic only (as long as they stay inside the grid's range).
\end{solution}

\begin{solution}{exo:rc:build-a-risk-engine:3}
VaR is a quantile, not a sum: a desk's VaR cannot be computed from its books' VaRs, nor the firm's from its
desks'. Keeping each trade's P\&L per scenario lets the engine compute any node's measure, and any new
grouping, by summing columns first.
\end{solution}

\begin{solution}{exo:rc:build-a-risk-engine:4}
The stand-alone ES of rates (47.34 million) and FX (14.63 million) add up to USD~61.96 million, from the unrounded values, against 50.89 million for the firm: 11.07
million of diversification. The Euler split puts almost all of the firm's ES on rates.
\end{solution}

\begin{solution}{exo:rc:build-a-risk-engine:5}
The desk is short options about a month from expiry. Over a 2.8\% move, gamma falls as the options move away
from the money, so losses grow more slowly than the parabola that extrapolates today's gamma.
\end{solution}

\begin{solution}{exo:rc:build-a-risk-engine:6}
The swaption's value depends on the forward swap rate, a weighted sum of \omterm{def:rc:curve-construction:pillar}{pillar} moves, and its convexity is in
the square of that sum: the cross products between \omterm{def:rc:curve-construction:pillar}{pillars}. Per-factor methods keep only the squares.
Cheap fixes: a grid on the trade's forward swap rate (one factor that carries the cross terms), a
parallel-shift grid plus \omterm{def:rc:curve-construction:pillar}{pillar} deltas for the rest, or \omterm{def:rc:build-a-risk-engine:reval}{full revaluation} of the few trades concerned.
\end{solution}

\begin{solution}{exo:rc:build-a-risk-engine:7}
Rate options in full: $200\times251 = 50\,200$; long-dated FX options by delta--gamma: $111\times19 = 2\,109$; swaps
by grid: $600\times49 = 29\,400$; short-dated FX options by grid: $89\times55 = 4\,895$. Total 86\,604, the count that the
engine's call counter returns for the plan.
\end{solution}

\begin{solution}{exo:rc:build-a-risk-engine:8}
The error depends on the book and the scenarios: new short-dated options, a longer horizon or a stressed
period can break an approximation that passed. The test must be re-run whenever either changes, and
routinely.
\end{solution}

\begin{solution}{pb:rc:build-a-risk-engine:1}
\textbf{1.} Full 251\,000, grid 50\,200, delta--gamma 17\,400.
\textbf{2.} Calls $\times\,50\times0.02/16/3\,600$: 4.36, 0.87 and 0.30 hours.
\textbf{3.} Grid and delta--gamma; \omterm{def:rc:build-a-risk-engine:reval}{full revaluation} misses a four-hour window.
\textbf{4.} The grid on trades, factors and points; \omterm{def:rc:build-a-risk-engine:reval}{full revaluation} on trades times scenarios.
\textbf{5.} $4\times3\,600\times16/(50\,000\times0.02) = 230.4$ calls per trade: 229 scenarios.
\textbf{6.} 1.59\% for the grid, 0.74\% for delta--gamma.
\textbf{7.} 8.64\% for the grid, 18.45\% for delta--gamma.
\textbf{8.} The short-dated FX options: their gamma changes over a ten-day move.
\textbf{9.} The swaptions: their convexity is in cross products of \omterm{def:rc:curve-construction:pillar}{pillar} moves.
\textbf{10.} Moves grow with the square root of the horizon, and the neglected terms grow faster than the
retained ones.
\textbf{11.} Grid by default, \omterm{def:rc:build-a-risk-engine:reval}{full revaluation} for the rate options, delta--gamma for the long-dated FX options.
\textbf{12.} 86\,604 calls, 1.50 hours.
\textbf{13.} 0.76\%, on the short-dated FX desk.
\textbf{14.} Yes: its delta--gamma error is 0.43\%, since options months from expiry have slowly changing gamma;
the grid's error there is larger (2.17\%) because seven points interpolate a nearly flat P\&L poorly.
\textbf{15.} The node errors against \omterm{def:rc:build-a-risk-engine:reval}{full revaluation}, on a schedule and whenever the book or scenarios change.
\textbf{16.} The short-dated FX desk: ten-day VaR USD~14.58 million against a limit of 12 million.
\textbf{17.} That its tail is a rates tail: rates contribute 47.34 of the 50.89 million of ES.
\textbf{18.} A run that finishes in time even in stress, with numbers reconciled and aggregated on a largely
automated basis.
\textbf{19.} Named result: \emph{the six o'clock deadline}: \omterm{def:rc:build-a-risk-engine:reval}{full revaluation} 4.36 hours (exact), grid 0.87 hours
(8.64\% worst error), delta--gamma 0.30 hours (18.45\%), plan 1.50 hours (0.76\%): at the ten-day horizon the plan
is the cheapest method within the 2\% budget; at one day, delta--gamma would be.
\textbf{20.} It depends on the book, the horizon and the scenarios, all of which change.
\end{solution}

\begin{solution}{iq:rc:build-a-risk-engine:1}
Trade and market-data sourcing (snapshots, hierarchy), scenario generation, revaluation through the pricing
library, the P\&L array, aggregation, measures and limits, reporting and lineage; batch orchestration,
failures handled per trade, and reconciliation with the books.
\iqlookfor{the stages and the P\&L array as the central object.}
\end{solution}

\begin{solution}{iq:rc:build-a-risk-engine:2}
Measure where the time goes; parallelise over trades; use grids or sensitivities where their error is small
and \omterm{def:rc:build-a-risk-engine:reval}{full revaluation} where it is not; reuse unchanged trades' results; cache market-data builds; fail single
trades without stopping the run.
\iqlookfor{measurement, parallelism and a plan checked against \omterm{def:rc:build-a-risk-engine:reval}{full revaluation}.}
\end{solution}

\begin{solution}{iq:rc:build-a-risk-engine:3}
\omterm{def:rc:build-a-risk-engine:reval}{Full revaluation} is exact and costs a call per trade and scenario. Grids cost calls per factor and capture
single-factor nonlinearity. Sensitivities are cheapest and fail on large moves and cross effects. The choice
is measured, desk by desk.
\iqlookfor{cost, error and when each fails.}
\end{solution}

\begin{solution}{iq:rc:build-a-risk-engine:4}
VaR is a quantile of a sum, not a sum of quantiles, and it is not subadditive in general. Allocate by Euler
contributions (each desk's expected loss in the firm's tail scenarios for ES), which add up to the total.
\iqlookfor{non-additivity and \omterm{def:rc:the-valuation-adjustment-desk:euler}{Euler allocation}.}
\end{solution}

\begin{solution}{iq:rc:build-a-risk-engine:5}
Instruments against closed forms and limits; approximations against \omterm{def:rc:build-a-risk-engine:reval}{full revaluation} on small moves and on
real scenarios; aggregation on fixtures and identities (Euler sums, node totals); cross-language twins on shared
fixtures; regression runs on stored outputs; timing tests.
\iqlookfor{a chain of tests, from pricing to reports.}
\end{solution}

\begin{solution}{iq:rc:build-a-risk-engine:6}
Governance, data architecture, accurate, complete, timely and adaptable risk data, and accurate, clear,
comprehensive reports. Banks struggle because of legacy systems, fragmented data, underfunded programmes
and weak board attention, as the Basel Committee's 2023 progress report and the ECB's 2024 guide noted.
\iqlookfor{the principles and the practical obstacles.}
\end{solution}
